\documentclass[11pt]{article}
\usepackage{amsmath,amssymb}
\usepackage{physics}
\usepackage[margin=1in]{geometry}
\usepackage{enumitem}

\title{ATPL 2025 --- Exercises: State Preparation and the Input Problem}
\author{Lecture 1: Why do we need hybrid programming models?}
\date{}

\begin{document}
\maketitle

Background: the "Input Problem" slides from Lecture 1. For the standard
recursive state-preparation construction referenced below, see e.g.\
M{\o}tt{\"o}nen, Vartiainen, Bergholm \& Salomaa, \emph{Transformation of
quantum states using uniformly controlled rotations} (2004), or Shende,
Bullock \& Markov (2006).

\begin{enumerate}[label=\textbf{Exercise \arabic*.}, leftmargin=2.2cm]

\item \textbf{Basis-state preparation is cheap.}
Given a target computational basis state $\ket{b_1 b_2 \cdots b_n}$ with
$b_k \in \{0,1\}$, describe the circuit that prepares it from
$\ket{0\cdots0}$ using only $X$ gates.
\begin{enumerate}[label=(\alph*)]
\item What is the total gate count?
\item What is the circuit \emph{depth} (can the $X$ gates all be applied
in parallel, or must they be sequential)? Contrast this with the general
amplitude state-preparation problem discussed below.
\end{enumerate}

\item \textbf{Counting the cost of a controlled-rotation state-prep
tree.}
The standard recursive construction for an arbitrary real-amplitude
$n$-qubit state uses one ``uniformly controlled'' $R_Y$ layer (a
rotation on qubit $n$, controlled on all $2^{n-1}$ possible values of
qubits $1,\dots,n-1$) to peel off one qubit at a time, then recurses on
the remaining $(n-1)$-qubit state.
\begin{enumerate}[label=(\alph*)]
\item Write down the recurrence $T(n) = 2\,T(n-1) + c\cdot 2^{n-1}$ for
the number of controlled-rotation gates, with $T(1)=O(1)$, justifying
the two terms.
\item Solve the recurrence (e.g.\ by unrolling it, or via a generating
function / substitution $T(n)/2^n = T(n-1)/2^{n-1} + c/2$) to find
$T(n)$ as an explicit function of $n$.
\item Compare your answer to the $\mathcal O(4^n)$ figure quoted in the
lecture. Is $\mathcal O(4^n)$ a tight bound for this construction, or a
looser (but easier to state) bound? Where might a factor like $4^n$
legitimately arise instead (hint: think about complex amplitudes, or
about a generic $2^n\times2^n$ unitary rather than just a state).
\end{enumerate}

\item \textbf{Complex amplitudes.}
Repeat Exercise 2, but now allow \emph{complex} amplitudes $a_i \in
\mathbb C$ (so each ``peeling off a qubit'' step needs both an $R_Y$
magnitude rotation \emph{and} an $R_Z$ phase rotation per branch). Does
the asymptotic growth rate (as a function of $n$) change from the
real-amplitude case, or only the constant factor?

\item \textbf{QRAM and the ``log-depth mirage''.}
A common claim is: ``QRAM lets you load $N$ data points in $\mathcal
O(\log N)$ time.'' The bucket-brigade QRAM architecture routes an
address query $\ket i$ down a binary tree of $\mathcal O(N)$ routing
qubits to reach memory cell $i$ in circuit \emph{depth} $\mathcal
O(\log N)$.
\begin{enumerate}[label=(\alph*)]
\item Suppose the $N$ values already sit in the QRAM's memory cells as
fixed physical parameters (like a lookup table burned into hardware).
Explain why the query
$\ket i\ket 0 \mapsto \ket i \ket{x_i}$
can indeed be performed in $\mathcal O(\log N)$ depth.
\item Now suppose instead that you want to encode $N$ \emph{arbitrary}
real amplitudes $a_0,\dots,a_{N-1}$ (freshly computed classically, not
sitting in a pre-built table) into the state $\sum_i a_i \ket i$.
Explain why this "loading" step still requires performing $N$
\emph{genuinely different} single-qubit rotations $R_Y(\theta_i)$ --
one per memory cell -- and therefore costs $\Omega(N)$ gates overall,
regardless of the $\mathcal O(\log N)$ query \emph{depth} of the
underlying QRAM routing structure.
\item Reconcile (a) and (b): under what circumstances does QRAM give a
genuine complexity advantage, and when is the "$\mathcal O(\log N)$"
figure misleading about the \emph{total} cost of getting your data into
the machine?
\end{enumerate}

\item \textbf{Structured states are cheap.}
Give an example of an $n$-qubit state with (generically) all $2^n$
amplitudes nonzero, that can nonetheless be prepared in $\mathrm{poly}(n)$
gates. (Hint: consider the state produced by applying the Quantum
Fourier Transform to a computational basis state, or a simple product
state written in a different basis, or the GHZ or $W$ states.) Contrast
this with a generic (Haar-random) state, which Exercises 2--3 show needs
exponentially many gates. What structural property distinguishes the
cheap states from the generic ones?

\end{enumerate}

\end{document}
